\chapter{Nephrostomy Placement}

\section*{Overview}
Nephrostomy placement inserts a tube through the skin into the kidney to drain urine when the ureter is blocked.

\section*{Alternative Names}
Percutaneous nephrostomy; nephrostomy tube insertion; PCN

\section*{Principle}
Under imaging guidance, a needle enters the renal pelvis and a catheter is placed to drain urine, bypassing an obstruction of the ureter.

\section*{History}
Percutaneous nephrostomy was developed in the 1950s and is now a routine interventional radiology procedure.

\section*{Why This Test or Procedure Is Ordered}
Performed for ureteral obstruction causing hydronephrosis, infection, or kidney injury, and to access the kidney for stone treatment.

\section*{Is This a Primary or Secondary Test or Procedure}
Primary drainage procedure for obstructive uropathy.

\section*{How This Test or Procedure Is Performed}
Performed under sedation with ultrasound and fluoroscopy. The collecting system is punctured, a wire is passed, and a drainage catheter is placed and secured.

\section*{What Preparation Is Needed}
Imaging, coagulation review, and antibiotics in selected cases.

\section*{Contraindications and Precautions}
Severe coagulopathy is a relative contraindication.

\section*{Risks}
Bleeding, infection or sepsis, injury to adjacent organs, and tube dislodgement or blockage.

\section*{What Happens After the Test or Procedure}
The tube is monitored and flushed; urine output is recorded.

\section*{Understanding Your Results}
Drainage and imaging confirm decompression of the kidney.

\section*{Normal Results}
Free drainage of urine with resolution of hydronephrosis.

\section*{Follow-up Tests and Procedures}
Tube exchange and imaging; definitive treatment of the obstruction.

\section*{References}
\begin{enumerate}
  \item MedlinePlus. Percutaneous nephrostomy. U.S. National Library of Medicine; 2025.
  \item Current Medical Diagnosis and Treatment. McGraw-Hill; 2025.
\end{enumerate}

\clearpage
